% =========================================================
% frontmatter/abstract.tex
% =========================================================

\chapter*{Abstract}
\thispagestyle{fancy}
\addcontentsline{toc}{chapter}{Abstract}
\begingroup
\setlength{\emergencystretch}{2em}

This report presents the analytical design and CST full-wave simulation of a 5G microstrip patch
antenna and a two-element phased patch array operating near 2.501~GHz. The design uses a
DiClad~880-IM substrate with $\varepsilon_r=2.20$, $\tan\delta=0.0009$, and
$h=1.524$~mm. Initial patch, feed-line, and inset dimensions were obtained from the
rectangular-patch transmission-line model and implemented as a parameterized CST model.

The final tuned single patch used a patch length of $38.40$~mm and an inset depth of $14.11$~mm.
Its simulated resonance occurred at $2.50507$~GHz with $S_{11,\min}\approx-39.87$~dB. The
PEC baseline produced peak directivity of approximately $7.226$~dBi, radiation efficiency of
$96.3\%$, and total efficiency of $96.2\%$. A supplementary fixed-mesh copper model preserved
the resonance, input match, directivity, and normalized pattern while reducing the radiation and
total efficiencies to $92.46\%$ and $92.38\%$, respectively. This comparison quantified the
conductor-loss effect omitted by the PEC baseline.

The element was extended to a two-port $2\times1$ array with center-to-center spacing
$d=59.935$~mm. Because mutual coupling changed the isolated-element resonant behavior, the array
required independent retuning. The final array used a patch length of $36.70$~mm and an inset
depth of $13.35$~mm. Near the operating frequency, the simulated results were approximately
$S_{11}=-11.93$~dB, $S_{22}=-15.20$~dB, and mutual coupling of about $-15$~dB.

Beam steering was demonstrated by combining the two CST port farfields with equal amplitudes and
relative phases of $0^\circ$, $-90^\circ$, and $+90^\circ$. The in-phase case produced a broadside
pattern with displayed peak directivity of approximately $9.655$~dBi. The phase-shifted cases
produced simulated beam displacements of approximately $20^\circ$ in opposite half-planes,
compared with the ideal two-element array-factor prediction of $30^\circ$ magnitude.

The study links closed-form synthesis, full-wave modeling, element tuning, conductor-loss
sensitivity, two-port characterization, mutual-coupling interpretation, and phase-controlled beam
steering. The results are analytical or simulated; no fabricated prototype or measured data are
claimed. Formal mesh- and boundary-convergence studies were not completed because of the CST
Learning Edition resource limit.
\par\endgroup
